\documentclass[../main.tex]{subfiles}

\begin{document}

\chapter*{Kata Pengantar}
\addcontentsline{toc}{chapter}{Kata Pengantar}

Puji syukur penulis panjatkan ke hadirat Tuhan Yang Maha Esa atas terselesaikannya buku ajar \textbf{Kriptografi} ini. Buku ini disusun sebagai panduan belajar bagi mahasiswa Program Studi Teknik Informatika untuk memahami prinsip-prinsip keamanan informasi melalui pendekatan kriptografi.

Buku ini dirancang dengan menggunakan kerangka \textit{Outcome-Based Education} (OBE), di mana setiap bab disusun untuk mencapai Capaian Pembelajaran Mata Kuliah (CPMK) dan Sub-CPMK yang terukur. Penulis menyadari bahwa kriptografi bukan sekadar menghafal algoritma, melainkan memahami landasan matematis dan logika di baliknya untuk dapat mengidentifikasi serta memitigasi ancaman keamanan di dunia nyata.

Materi dalam buku ini mencakup evolusi kriptografi, mulai dari cipher klasik, kriptografi simetris modern (DES, AES), kriptografi asimetris (RSA, Diffie-Hellman, ElGamal, ECC), hingga mekanisme integritas data menggunakan fungsi hash dan tanda tangan digital. Sebagian besar materi disarikan dari buku \citet{munir2019} dan dilengkapi dengan rujukan standar di bidang kriptografi.

Penulis menyadari bahwa buku ini masih jauh dari sempurna, oleh karena itu kritik dan saran yang membangun sangat diharapkan demi perbaikan di masa mendatang. Ucapan terima kasih kepada pihak-pihak yang telah mendukung penyusunan buku ini disampaikan pada halaman \emph{Ucapan Terima Kasih}.

\vspace{2cm}

\begin{flushright}
Bandung, 2026\\
\vspace{1cm}
\textbf{Rinaldi Munir}
\end{flushright}

\end{document}
